\documentclass[12pt]{article}
\usepackage[margin=1in]{geometry}
\usepackage{amsmath,booktabs,graphicx,threeparttable}

\begin{document}

\providecommand{\StructuralRobustnessPath}{appendix/structural-robustness}

\subsection{Robustness to a Finite-Mixture Distribution of Intrinsic Motivation}
\label{app:finite-mixture-types}

The heavy-tail exercise retains a symmetric, unimodal distribution. We also
estimate a two-component Gaussian mixture for intrinsic motivation, allowing
the data to generate asymmetry, shoulders, or bimodality. To preserve the
scale normalization of the baseline model, the mixture is parameterized to
have mean zero and variance one. Its free shape parameters are the mass on the
lower component and the share of total variance attributable to separation
between the component means. We assign these parameters $\operatorname{Beta}
(2,2)$ and $\operatorname{Beta}(2,4)$ priors, respectively. The remaining
structural model and all other priors are unchanged.

The mixture fit uses four chains with 800 warmup and 800 retained draws per
chain. The maximum split $\widehat R$ is 1.004, the minimum bulk and tail
effective sample sizes are 852 and 1,180, respectively, and there are no
divergences or maximum-treedepth transitions. The posterior median mass on the
lower component is 0.61 (95 percent credible interval 0.12--0.92), and the
median between-component share of variance is 0.31 (0.05--0.68). The broad
shape intervals are appropriate here: the exercise deliberately allows a
richer latent distribution rather than presuming that two sharply separated
classes are present.

Figure~\ref{fig:finite-mixture-density} shows the estimated latent-type
density and its pointwise 95 percent posterior interval. The standard-normal
benchmark is shown for comparison. The posterior permits appreciable
non-Gaussian shape but places its central density close to the Gaussian
benchmark; thus the exercise is flexible without mechanically forcing a
second mode.

\begin{figure}[htbp]
  \centering
  \includegraphics[width=0.72\textwidth]{\StructuralRobustnessPath/figures/finite-mixture-density.pdf}
  \caption{Estimated finite-mixture distribution of intrinsic motivation}
  \label{fig:finite-mixture-density}
  \begin{minipage}{0.88\textwidth}\footnotesize
    \textit{Notes:} The solid line is the pointwise posterior median density
    and the shaded region is its pointwise 95 percent posterior interval. The
    dashed line is the standard-normal density used in the baseline model;
    dotted lines show the two weighted components evaluated at the posterior
    median mixture parameters. Every mixture draw is normalized to mean zero
    and variance one.
  \end{minipage}
\end{figure}

Table~\ref{tab:finite-mixture-multipliers} compares social multipliers from
the baseline and mixture specifications. The posterior medians and arm
ordering are similar at all three distances: multipliers are largest for
Control, remain above one for Calendar, and are attenuated below one for Ink
and Bracelet. At 1.5 km the 95 percent intervals exclude one for Control and
just exclude one for Bracelet, while those for Calendar and Ink include one;
at 2.5 km the intervals exclude one for Control, Ink, and Bracelet. The
mixture fit is less precise, especially for the local 500-meter derivative,
whose intervals can include negative values. We therefore interpret this
exercise as evidence that the median multiplier pattern is not an artifact of
imposing a single Gaussian latent type, while recognizing the additional
uncertainty introduced by flexibly estimating the latent distribution.

\begin{table}[htbp]
  \centering
  \small
  \caption{Social multipliers under Gaussian and finite-mixture intrinsic motivation}
  \label{tab:finite-mixture-multipliers}
  \begin{threeparttable}
    \begin{tabular}{lccc}
\toprule
Treatment & 500m & 1500m & 2500m \\
\midrule
\multicolumn{4}{l}{\textit{Panel A: Gaussian intrinsic motivation}} \\
Bracelet & 0.89 & 0.83 & 0.78 \\
 & (0.70, 1.04) & (0.69, 0.95) & (0.67, 0.90) \\ 
Calendar & 1.26 & 1.22 & 1.21 \\
 & (1.11, 1.47) & (1.05, 1.46) & (1.01, 1.52) \\ 
Ink & 0.92 & 0.86 & 0.81 \\
 & (0.80, 1.03) & (0.74, 0.97) & (0.70, 0.94) \\ 
Control & 1.43 & 1.48 & 1.57 \\
 & (1.24, 1.73) & (1.24, 1.87) & (1.26, 2.03) \\ 
\addlinespace
\multicolumn{4}{l}{\textit{Panel B: Two-component Gaussian-mixture intrinsic motivation}} \\
Bracelet & 0.84 & 0.81 & 0.77 \\
 & (-0.04, 1.07) & (0.56, 1.00) & (0.53, 0.94) \\ 
Calendar & 1.21 & 1.19 & 1.20 \\
 & (-0.08, 1.54) & (0.85, 1.50) & (0.88, 1.55) \\ 
Ink & 0.90 & 0.85 & 0.81 \\
 & (0.63, 1.08) & (0.60, 1.01) & (0.58, 0.98) \\ 
Control & 1.38 & 1.46 & 1.56 \\
 & (-0.01, 1.82) & (1.03, 1.91) & (1.19, 2.03) \\ 
\bottomrule
\end{tabular}

    \begin{tablenotes}\footnotesize
      \item \textit{Notes:} Entries are posterior medians of the point social
      multiplier evaluated at the indicated distance. Parentheses contain 95
      percent posterior credible intervals. The mixture has two Gaussian
      components and is constrained to have mean zero and variance one.
    \end{tablenotes}
  \end{threeparttable}
\end{table}

\clearpage

\end{document}
